\section{Introduction}
\label{sec:intro}

In plain terms, this paper asks whether any nearby star is transmitting a
narrow radio signal that ALMA happened to record while observing
something else. ALMA's public archive holds thousands of hours of
millimetre and submillimetre observations taken for quite different
reasons, mostly studies of dust discs and planet formation. Every one of
them also contains a spectrum of the star at the centre of the field. A
transmitter would appear in such a spectrum as a spike confined to one
or two channels. It would sit at the star's position and nowhere else in
the field, and reappear if the star were observed again. We searched the
archive for exactly that, with one pipeline and no human intervention. Nothing of the kind survived.

A technological civilisation may imprint the electromagnetic spectrum
detectably whether or not it intends to, a \emph{technosignature}
\citep{Sheikh2020,SETIreview2025}. Millimetre wavelengths are an
attractive place to look for one. Interstellar scattering and dispersion
fall steeply with frequency, so a high-frequency carrier arrives less
distorted. A given aperture delivers higher gain, and so more radiated
power per watt supplied to the transmitter. Coherent millimetre emission
is also already engineered on Earth, for directional links, radar and
beamed power.

The band is correspondingly hard to search. The atmosphere transmits
unevenly and adds phase noise, and receivers are noisier than at
centimetre wavelengths. Archived correlator channels are thousands of
times coarser than a dedicated narrowband backend, so a carrier is
diluted within one channel and its drift can smear it across several.
These bands are also crowded with molecular lines, a dense background of
real astrophysical signals that a threshold search must separate from an
artificial one.

Six decades of radio searches have examined many thousands of stars, yet
sample only \HaystackFrac{} of the space of transmitter frequency,
bandwidth, power and duty cycle \citep{Wright2018}. Essentially all of
that effort lies below ${\sim}30$\,GHz; above it we are aware of one
published search, in ALMA Band~3 \citep{Mason2025}. Few searches at any
frequency reprocess existing interferometric archives, although those
archives already hold a spectrum of every target they have pointed at
(\S\ref{sec:background}).

This paper is an archival ALMA search of nearby stellar systems across
that band, with its sensitivity and its selection measured rather than
assumed. The signal class is an \emph{unresolved spectral carrier}:
emission confined to a few channels, unresolved at the stellar position,
and persistent enough to recur between epochs. Each qualifier is a real
restriction. A carrier broader than the channel width would not be
recovered here, nor one that fills the synthesised beam, nor one that
appears once and never again. The fine channels used are
\ChanLoA--\ChanHiA\,kHz wide, so ``narrowband'' here describes the
assumed transmitter and never the resolving power of the data.
Channelisation also divides the archive. Only \NWinA{} of the
\NWindows{} searched windows support the fine-channel search; the rest
register an unexplained spike but cannot tell a drifting signal from a
stationary one.

Every window follows the same four steps: a threshold crossing;
attribution to a known line; comparison of the star with \NCtrl{}
control positions in the same field; and a search for recurrence at
another epoch. Figure~\ref{fig:funnel} is the roadmap. One
worked case runs through the paper: a crossing toward $\beta$~Pictoris,
the largest star-versus-control contrast here, attributed to the disc's
known circumstellar CO in two transitions and \NStageOneBpicEb{}
execution blocks (\S\ref{sec:technosearch}). Where this paper says
``sensitivity'' it means $\mathrm{EIRP}_{90}$, the equivalent isotropic
radiated power that the complete automated selection recovers nine times
in ten. The internal trigger level $P_{\rm trig}$ is not a sensitivity.

Three limits on interpretation hold throughout. The targets were chosen
by archival disc and planet-formation programmes, not for this search.
The sample is therefore enriched in F and G stars and depleted in M
dwarfs: only \NMSearched{} of \NMCensus{} catalogued M~dwarfs within
40\,pc appear in it. Every statistical statement here concerns the ALMA
archive-selected stellar sample, not the solar neighbourhood. The survey
cannot constrain the local habitable-zone planet population, which
M~dwarfs dominate (\S\ref{sec:bias}). The search statistic and the
line mask were defined on these data, so the rates quoted here are post
hoc, and \S\ref{sec:method} sets out what was fixed when. And only
\EpIndep{} of \EpTot{} systems were observed more than a day apart, so a
non-detection constrains a persistent carrier far better than an
intermittent one.

%% REQUEST (appendix-triage owner): the plain-language opening now lives
%% here, in the first paragraph. Please delete it from
%% sections/app_00_overview.tex so it is not stated twice.
%% REQUEST (background owner): three statements have moved up into this
%% section and should now be cut from sections/02_background.tex, where
%% they are duplicated -- (i) "There are physical reasons to go higher:
%% interstellar scattering ... without forbidding it"; (ii) "Essentially
%% all of that work lies below ~30 GHz, above which only Mason2025 has
%% searched, in ALMA Band 3"; (iii) the haystack-fraction sentence if it
%% recurs there. Section 2 should begin with the survey-by-survey
%% review, not with the motivation.
%% REQUEST (figures owner): \ref{fig:funnel} must survive as a body
%% figure, or this paragraph loses its roadmap pointer.
